\documentclass[aspectratio=169,11pt]{beamer}

% ============================================
% PACKAGES
% ============================================
\usepackage[utf8]{inputenc}
\usepackage[T1]{fontenc}
\usepackage[ngerman]{babel}
\usepackage{lmodern}
\usepackage{tcolorbox}
\tcbuselibrary{skins, hooks}
\usepackage{listings}
\usepackage{xcolor}
\usepackage{fontawesome5}
\usepackage{tikz}
\usetikzlibrary{trees, shapes, arrows.meta}

% Modern Font
\IfFileExists{FiraSans.sty}{\usepackage[sfdefault]{FiraSans}}{}
\renewcommand{\familydefault}{\sfdefault}

% ============================================
% COLORS & DESIGN SYSTEM
% ============================================
\definecolor{BgColor}{HTML}{FAFAFA}
\definecolor{Primary}{HTML}{2B3A42}
\definecolor{Accent}{HTML}{E84A5F}
\definecolor{Secondary}{HTML}{3F5765}
\definecolor{CodeBg}{HTML}{F0F0F0}
\definecolor{Success}{HTML}{28A745}

\setbeamercolor{background canvas}{bg=BgColor}
\setbeamercolor{title}{fg=Primary}
\setbeamercolor{frametitle}{bg=Primary, fg=white}
\setbeamercolor{structure}{fg=Accent}
\setbeamercolor{normal text}{fg=Primary}
\setbeamercolor{itemize item}{fg=Accent}

\setbeamertemplate{navigation symbols}{}
\setbeamertemplate{footline}{%
  \leavevmode%
  \hbox{%
  \begin{beamercolorbox}[wd=.333333\paperwidth,ht=2.25ex,dp=1ex,center]{author in head/foot}%
    \usebeamerfont{author in head/foot}\tiny Falko, Malte, Max, Jana
  \end{beamercolorbox}%
  \begin{beamercolorbox}[wd=.333333\paperwidth,ht=2.25ex,dp=1ex,center]{title in head/foot}%
    \usebeamerfont{title in head/foot}\tiny Grundlagen der Praktischen Informatik
  \end{beamercolorbox}%
  \begin{beamercolorbox}[wd=.333333\paperwidth,ht=2.25ex,dp=1ex,right]{date in head/foot}%
    \usebeamerfont{date in head/foot}\insertframenumber{} / \inserttotalframenumber\hspace*{2ex} 
  \end{beamercolorbox}}%
  \vskip0pt%
}

% ============================================
% CUSTOM BOXES
% ============================================
\tcbset{
    modernbox/.style={
        enhanced,
        boxrule=0pt,
        frame hidden,
        colback=white,
        coltitle=white,
        colbacktitle=Secondary,
        fonttitle=\bfseries,
        drop lifted shadow,
        arc=4pt,
        toptitle=2mm,
        bottomtitle=2mm,
        left=2mm,
        right=2mm
    },
    alertbox/.style={
        modernbox,
        colbacktitle=Accent
    }
}

% ============================================
% CODE LISTINGS
% ============================================
\lstdefinelanguage{Haskell}{
  keywords={class, data, deriving, do, else, if, in, import, let, module, of, then, type, where, case, error},
  keywordstyle=\color{Accent}\bfseries,
  identifierstyle=\color{Primary},
  sensitive=true,
  comment=[l]{--},
  commentstyle=\color{gray}\ttfamily\itshape,
  stringstyle=\color{teal}\ttfamily,
  morestring=[b]"
}

\lstset{
    language=Haskell,
    backgroundcolor=\color{CodeBg},
    basicstyle=\ttfamily\scriptsize,
    breaklines=true,
    frame=single,
    rulecolor=\color{CodeBg},
    frameround=tttt,
    numbers=left,
    numberstyle=\tiny\color{gray},
    xleftmargin=15pt,
    tabsize=4,
    showstringspaces=false
}

% ============================================
% TITLE PAGE
% ============================================
\title{\textbf{Tutorium 7}}
\subtitle{Bäume \& Altklausuraufgaben}
\institute{Grundlagen der Praktischen Informatik}
\author{}
\date{\today}

\begin{document}

\begin{frame}[plain]
    \titlepage
\end{frame}

\begin{frame}{Inhaltsverzeichnis}
    \tableofcontents
\end{frame}

% ============================================
% BÄUME
% ============================================
\section{Bäume}

\begin{frame}[fragile]{Binärbäume in Haskell}
    Bäume sind rekursive Datenstrukturen. Ein Baum kann leer sein (\texttt{Nil}) oder ein Knoten (\texttt{Node}) mit einem Wert und zwei Teilbäumen (links und rechts).
    
    \begin{tcolorbox}[modernbox, title=Definition eines Binärbaums]
\begin{lstlisting}
data BinTree a = Nil | Node a (BinTree a) (BinTree a)
  deriving (Show, Eq)
\end{lstlisting}
    \end{tcolorbox}
    \small
    \textbf{Beispiel eines Knotens (Blatt):} \texttt{Node 3 Nil Nil}

    \vspace{0.3cm}
    \begin{center}
    \begin{tikzpicture}[level distance=0.8cm, sibling distance=1.5cm, every node/.style={circle, draw, minimum size=0.5cm}]
    \node {8}
      child {node {3}
        child {node {1}}
        child {node {6}}
      }
      child {node {10}
        child {node[draw=none] {}}
        child {node {14}}
      };
    \end{tikzpicture}
    \end{center}

\end{frame}

\begin{frame}[fragile]{Operationen auf Bäumen}
    Analog zu Listen werden Operationen auf Bäumen über \textbf{Pattern Matching} realisiert. Wir definieren einen Fall für \texttt{Nil} und einen für \texttt{Node}.
    
    \begin{columns}[T]
        \begin{column}{0.48\textwidth}
            \begin{tcolorbox}[modernbox, title=Knotensumme]
\begin{lstlisting}
sumTree :: Num a => BinTree a -> a
sumTree Nil = 0
sumTree (Node v l r) = 
    v + sumTree l + sumTree r
\end{lstlisting}
            \end{tcolorbox}
        \end{column}
        \begin{column}{0.48\textwidth}
            \begin{tcolorbox}[modernbox, title=Knoten zählen]
\begin{lstlisting}
countNodes :: BinTree a -> Int
countNodes Nil = 0
countNodes (Node _ l r) = 
    1 + countNodes l + countNodes r
\end{lstlisting}
            \end{tcolorbox}
        \end{column}
    \end{columns}
\end{frame}

\begin{frame}[fragile]{Tree Fold (Fold für Bäume)}
    Da die rekursiven Traversierungen immer ähnlich aussehen, können wir eine höhere Funktion \texttt{foldTree} definieren:
    
    \begin{tcolorbox}[modernbox, title=\texttt{foldTree}]
\begin{lstlisting}
foldTree :: (a -> b -> b -> b) -> b -> BinTree a -> b
foldTree _ z Nil = z
foldTree f z (Node v left right) = 
    f v (foldTree f z left) (foldTree f z right)
\end{lstlisting}
    \end{tcolorbox}
    
    Damit lassen sich \texttt{sumTree} und \texttt{countNodes} extrem kurz schreiben:
\begin{lstlisting}
sumTree' = foldTree (\v l r -> v + l + r) 0
countNodes' = foldTree (\_ l r -> 1 + l + r) 0
\end{lstlisting}
\end{frame}

% ============================================
% ALTKLAUSUR
% ============================================
\section{Altklausur \& Blatt 4}

\begin{frame}[fragile]{Aufgabe: Elemente Zählen (mit Filter vs. Fold)}
    \textbf{Klausuraufgabe:} Implementieren Sie \texttt{count}, die die Anzahl der Elemente zurückgibt, die ein Prädikat erfüllen.
    
    \pause
    \begin{tcolorbox}[modernbox, title=Lösungswege]
\begin{lstlisting}
-- Mit Filter (einfach!)
count' :: (a -> Bool) -> [a] -> Int
count' p xs = length (filter p xs)

-- Mit Fold
count'' :: (a -> Bool) -> [a] -> Int
count'' p xs = foldr (\x acc -> if p x then acc + 1 else acc) 0 xs
\end{lstlisting}
    \end{tcolorbox}
\end{frame}

% ============================================
% LIVE DEMO
% ============================================
\section{Praxis}

\begin{frame}[plain]
    \begin{center}
        \Huge \textbf{Live Demo} \\
        \vspace{0.5cm}
        \large \faLaptopCode~ \texttt{01\_altklausuraufgaben.hs}, \texttt{02\_baeume.hs}, \texttt{sheet04.hs} \\
        \vspace{0.3cm}
        \normalsize Wir programmieren Bäume und besprechen Blatt 4 sowie weitere Klausuraufgaben!
    \end{center}
\end{frame}

\end{document}
